\chapter{Security engineering, privacy, and threat modelling}

\section{Security is a system property}
Security engineering protects assets against unwanted actions under a stated
threat model. It includes confidentiality, integrity, availability, authenticity,
accountability, and privacy. A library patch cannot compensate for an endpoint
that authorizes the wrong tenant or logs credentials.

Start by naming assets and consequences:

\noindent\begin{tabularx}{\textwidth}{@{}p{0.23\textwidth}p{0.26\textwidth}X@{}}
\toprule
Asset & Adversary or accident & Consequence \\
\midrule
Order data & Cross-tenant read or export error & Confidentiality and legal exposure \\
Payment state & Forged callback or replay & Financial loss and reconciliation failure \\
Credentials & Log leak or phishing & Account takeover \\
Deployment pipeline & Malicious dependency or stolen token & System-wide compromise \\
\bottomrule
\end{tabularx}

\section{Threat modelling}
A lightweight threat model records actors, assets, entry points, trust boundaries,
and abuse cases. STRIDE is one categorization technique; attack trees, misuse
cases, and data-flow diagrams are alternatives. The technique is useful when it
changes a design or test. Listing threats without owners or mitigations is a
worksheet, not security engineering.

For an order API, ask:

\begin{enumerate}
  \item Can an unauthenticated caller reach it?
  \item How is tenant ownership derived and checked on every read and write?
  \item Can a caller replay a payment callback or idempotency key?
  \item Are quantities, IDs, URLs, files, and error messages bounded and encoded?
  \item Which secrets cross the process, logs, build system, and backups?
\end{enumerate}

\section{Secure defaults and input handling}
Default-deny authorization, least privilege, parameterized database queries,
context-appropriate output encoding, safe error messages, bounded input, and
secure transport are durable principles. “Validate input” is incomplete: validate
syntax at the boundary, enforce domain rules in the domain, and enforce critical
invariants in storage. Canonicalize before comparing where multiple encodings
could represent the same resource.

Authentication identifies a principal; authorization decides whether that
principal may perform a specific action on a specific resource. An authenticated
user is not automatically authorized to read every order. Check authorization
server-side, close to the resource access, and test both allowed and denied paths.

The OWASP ASVS provides a versioned basis for testing web application security
controls and requirements; its identifiers can change between versions, so cite
the exact version when using a requirement \cite{owasp2025asvs}.

\section{Secrets and cryptography}
Do not store secrets in source control, client bundles, ordinary logs, or error
messages. Use a secret manager or protected environment injection, rotate keys,
limit access, and record a plan for compromise. Choose well-reviewed cryptographic
primitives through a maintained library. Do not invent encryption or password
hashing schemes. Encryption does not solve authorization, key management, or
metadata leakage.

\section{Privacy by design}
Privacy includes purpose limitation, data minimization, access control,
retention, deletion, transparency, and accountability. A personal-data field
should have a reason, an owner, a retention rule, and an access path. Test exports
and deletion across caches, replicas, logs, backups, analytics, and derived data.
Legal obligations vary by jurisdiction; this chapter offers engineering
controls, not legal advice.

\section{Supply-chain security}
Record dependency versions and integrity, review transitive risk, protect build
credentials, produce provenance where possible, and make the release artifact
reproducible enough to investigate. NIST's SSDF provides a high-level vocabulary
for integrating secure development practices into a lifecycle \cite{nist2022ssdf}.
As of this edition's research date, NIST had also published an initial public
draft for a revision; a draft should not be represented as a final standard
\cite{nist2025ssdfdraft}.

\section{Security verification}
Use layered evidence: code review, static checks, dependency review, unit and
integration tests, dynamic testing, configuration checks, secret scanning, and
authorized penetration testing. Keep a vulnerability response path with
severity, owner, affected versions, remediation, disclosure, and verification.
No checklist proves absence of vulnerabilities. The aim is to reduce likely
failure, limit impact, and improve recovery.

\begin{exercise}
Draw a data-flow diagram for a multi-tenant order API. Mark trust boundaries and
write five abuse cases. For each, choose one preventive control and one test that
would detect a regression.
\end{exercise}

\section{Chapter summary}
Security is a property of architecture, code, data, people, and operations.
Model assets and adversaries, separate authentication from authorization, use
secure defaults and layered validation, minimize and govern personal data, and
keep standards versioned and evidence-based.
